\section{Completion of the local estimate}
\label{sec:completion}

\subsection{Restarts on the original summands}

We complete the target for a base segment \(R_0\) by the
construction in \cite[Section~10, One attempt and its restarts]{H}.
On an attempt \(R\), the main realization and actual uses precede
the prospective-offset calculation; the chosen whole sides then
precede helper closure and its dependent size and counting cuts.
The parent functions and operators remain fixed when original
summands are transferred at a cut.

The cuts accumulate nonnegative partial path counts, component
square sizes, or partial ancestor suprema.  Their integrated
bounds permit the thresholds to be increased by fixed factors.
Together with the initial marking-failure cuts, the union of first
stopped cells can then be chosen to have total length at most \(|R|/4\).
For the strong stack thresholds, only auxiliary choices contribute
a cardinality overhead, absorbed in \(G_0\); all full indices
\(\nu=(G,b,\chi)\) are controlled jointly by the aggregate energy
and full-hull bounds.  The polynomially many configurations in
\(P,k\) occur in the power budgets.  The hull construction
includes every potentially needed ancestor within a block before
the common cuts, and its maximal functions remain defined on
whole retained hull cells, including portions below a cut.
Large-frozen-complement flags are charged to the powered error
sum; only a first crossing of their accumulated budget creates
a restart root.

At a stopped cell \(R'\), send the original base summands there
and at its descendants to a fresh attempt.  Their original
depths, input schedules, and basic expressions remain assigned
to them.  The restrictions \(F_{j,s}|_{R'}\) satisfy the same
block normalization in each slot.  On \(R'\), the input size bound
\eqref{eq:h-input-size} gives the ordinary \(L^2\) budget
\(|R'|^{1/2}\), and \eqref{eq:h-input-variation} has right side
\(5a|R'|\).  Rebase the inherited basic columns with their
private constants unchanged and use genuine local refits and
inherited continuations on the new root.  The old continuations
needed by the parent's summands remain fixed.

All root bounds are obtained anew on \(R'\).
Equation~\eqref{eq:h-projection-energy} gives the basic rolling
energy with factor \(|R'|\).  The proofs of
\eqref{eq:h-certified-variation} and
\eqref{eq:h-frozen-composition} apply to the new main and mode
constructions, and their newly constructed linear sides satisfy
\eqref{eq:h-linear-energy} on \(R'\).  In particular, this
linear energy is proved by the new-root construction, with no
inference of proportional child energy from a parent's arbitrary
family.  The local sizes, column budgets, and inherited label
properties used in those proofs hold on \(R'\) as well.  The displayed
cross-depth bounds and local sizes thus have the same constants on
every new root, independently of the remaining number of depths.

The packing-of-attempts Lemma of \cite[Section~3]{H} now gives
\[
 \sum_{\substack{R\ {\rm attempt\ root}\\
                  {\rm for\ the\ segment}\ R_0}}|R|
 \le |R_0|\sum_{n\ge0}4^{-n}=\tfrac43|R_0|.
\]
The recursion is finite on the original finite depth tree;
equivalently, induction on the remaining depths gives this
geometric bound.

\subsection{The estimate at one accuracy}

At each \(I\in\mathcal A_R^{\rm ret}\), apply the minimum to the
absolute value of its whole contribution.  Use
\eqref{eq:h-localization} and subadditivity to separate the fixed
search leaves while keeping the group sum inside the cell.
In each leaf expand \eqref{eq:h-core-side}.  Terms with at least
three cores use \eqref{eq:h-core-estimate}; terms with at least
three sides use the \(p\)th-power estimate of the side calculation,
where \(p=1-\varepsilon>2/3\) is supplied by the core calculation.
Exactly two sides use the transferred main term and its freeze,
large-complement, and restoration estimates.  These cases exhaust
the four-position expansion.

For each powered within-cell contribution use
\[
 \min(\eta,x)\le\eta^{1-p}x^p\qquad(x\ge0).
\]
The near terms use \eqref{eq:h-near-band} at power one, with
their coupled band sums kept inside the cell.  Physical and
numerical errors are also estimated at power one.  Increase
\(C_4\) to include the powered costs and the smaller main cost
\(\eta P^C|R|\).  The overhead from source leaves and auxiliary
choices fits \(G_0\) for the strong terms; the polynomially many
power-budget choices fit the enlarged \(C_4\).  The strong energy
and hull bounds already sum all groups and blocks.  Thus on the
retained cells of one attempt,
\begin{equation}
 \begin{aligned}
 \sum_{I\in\mathcal A_R^{\rm ret}}|I|
        \min(\eta,|H_I(z_I)|)
 &\le C_q\eta^\varepsilon P^{C_4}|R|\\
 &\quad +(1+C_N(q))C_qP^{-1/2}G_0
                (K+2)^{C_{19}(q-2)}|R|\\
 &\quad +O(e^{-k}|R|).
 \end{aligned}
 \label{eq:h-retained-estimate}
\end{equation}
This is the conditional retained-cell calculation of
\cite[Section~10]{H}, with its basic square, variation, and
frozen-row hypotheses supplied above.  Summing it over the
attempt roots by the preceding packing bounds the full segment
target on \(\mathcal A_{R_0}\) by the same right-hand expression
with \(|R_0|\), up to a fixed factor.

\subsection{\texorpdfstring{Choice of \(q\), \(\alpha\), and starting accuracy}{Choice of q, alpha, and starting accuracy}}

The degrees in \eqref{eq:h-retained-estimate} and
\eqref{eq:h-budget-hierarchy} can be fixed independently of
\(q\).  The main-row and stack calculation following
\eqref{eq:h-stacks} records the norm and squared losses:
row logarithms, Compression of \(O(\log k)\) main factors,
and the strong stack factors fit \(G_0\); the larger \(D_*\)
is confined to the mode construction, Bessel, and counting estimates.  The base
row, smooth, and delay estimates add fixed polynomial degrees
in \(k\).  Powers of \(K\) enter the allowed component and top
budgets and the displayed short-block factor.  These facts
preserve absolute upper bounds on \(C_4,C_5,C_{19}\).

We need the fixed-input constant at the same \(q\).
The \emph{Uniform local estimate} of \cite[Section~2]{H} is
stated for some \(q\in(2,3)\), and its proof permits the choice
required here.  The \emph{Local detection} Lemma of
\cite[Section~4]{H} holds for every \(q>2\), with an absolute
logarithmic degree and \(q\)-dependence only in its coefficient.
The \emph{Short-block comparison} Lemma of
\cite[Section~7]{H} holds for every \(2<q<3\), with the explicit
power \(11(q-2)/(q-1)\) used above.  The other degree bounds in
H's hierarchy are fixed before \(q\).  Its
\emph{Order of the parameters} in \cite[Section~10]{H} requires
\[
 C_{19}^{\rm H}C_5^{\rm H}(q-2)<\tfrac18,
\]
where the superscript denotes the source degree bounds.
The source stagger and quality exponents precede \(q\), while
its power parameter and starting accuracy follow \(q\).

Choose the one \(q\in(2,3)\) for this proof so close to two that
\begin{equation}
 \max\{C_{19}^{\rm H}C_5^{\rm H},\,C_{19}C_5\}(q-2)<\tfrac18.
 \label{eq:h-common-q}
\end{equation}
Run the fixed-input parameter and absorption proof of
\cite[Sections~2 and~10]{H} at this chosen value, with its own
subsequent choices of power parameter and starting accuracy.
Its normalized absorption Proposition and outer reduction give
\begin{equation}
 C_{\rm H}(q):=\sup_{N\ge1}C_N(q)<\infty
 \label{eq:h-uniform-fixed}
\end{equation}
for the fixed compactly supported \(L^q\) functions in
\eqref{eq:h-fixed-local}.  This is the consequence of H's
proof-level parameter order at the one selected \(q\).

Now choose our \(\alpha>0\) with
\(\alpha C_4<\varepsilon/2\); \(C_4\) already dominates the
powered side costs.  For this fixed \(\alpha\),
\(P=\lceil e^{\alpha k^\theta}\rceil\) dominates every fixed
power of \(k\) and \(G_0=P^{o(1)}\).  By
\eqref{eq:h-common-q}, after increasing the starting \(k\),
\[
 \begin{gathered}
 P^{-1/2}G_0(K+2)^{C_{19}(q-2)}\le C_qP^{-1/4},\\
 \eta^\varepsilon P^{C_4}\le C_q\eta^{\varepsilon/2},
 \qquad P^{-1/4}\le\eta^{\alpha/4}.
 \end{gathered}
\]
Indeed the \(K\) power costs less than \(P^{1/8}\), and
\(G_0\le P^{1/8}\) for large \(k\).  Since \(0<\theta<1\),
the numerical error \(e^{-k}\) is smaller than the resulting
quality power.  Substitute \eqref{eq:h-uniform-fixed} into
\eqref{eq:h-retained-estimate} and sum the attempt roots.
For a fixed \(c_*>0\), this gives the full segment bound
\begin{equation}
 \sum_{I\in\mathcal A_{R_0}}|I|
       \min(\eta,|H_I(z_I)|)
 \le C_q\eta^{c_*}|R_0|.
 \label{eq:h-quality-gain}
\end{equation}
The fixed constant \(C_{\rm H}(q)\) is included in \(C_q\).

\subsection{Summing the accuracy expansion}

Use the final accuracy summation of
\cite[Section~10, Summing the base expansion]{H}.  For a tuple
with later level \(t\), the earlier level has \(O(t)\) choices
and there are only a bounded number of slot patterns.  The
intersections of the tuple's bounded number of basic segment
families have total root length at most a fixed multiple of
\(|J|\); the localization attempt packing preserves that bound.
The noninitial costs in \eqref{eq:h-quality-gain} therefore
sum because
\[
 \sum_{t\ge12}t\,e^{-c_*k_t^\theta}<\infty,\qquad
 k_{t+1}=k_t^A,\quad A>1.
\]
The separate global replacement errors in
\eqref{eq:hbase-high-replacement} and the two-\(U\) errors have the
same summable multiplicity.  The
finitely many initial levels have the depth-independent bound
from the base active-scale reduction.

Finally take the two-earliest-increments expansion at a finite
accuracy cutoff.  At each fixed cell all its accuracy updates
were evaluated on the same actual \(F_{j,s(I)}\); the per-cell
pair comparisons in the base proof therefore make the terminal
residual tend to zero.  There are finitely many original cells,
so this limit commutes with their finite sum.  The triangle
inequality, the uniform summable bounds, and passage to that
cutoff limit give \(C|J|\) for the original sum.

The cross-depth estimates
\eqref{eq:h-input-variation}, \eqref{eq:h-row-use},
\eqref{eq:h-projection-energy}, and \eqref{eq:h-smooth-use},
together with the single-depth size bound
\eqref{eq:h-input-size}, have constants independent of the four
input block partitions and their numbers \(W_j\).  The resulting
constant is therefore uniform in the input schedules, as
required by Theorem~\ref{thm:depth-local}.
